% =============================================================
%  Subsection — Natural power (surge-impedance loading)
%  Part of design/ac-transmission; \input from ac-transmission.tex.
%
%  Figures in ./fig (np_*.pdf), all vector, generated by
%  fig/natural_power_figures.py from the exact two-port equations
%  (numpy + matplotlib only).  Cross-checked against a 20-segment
%  pandapower model; agreement better than 0.1 % on every quantity.
%
%  Reference line throughout: 400 kV, 2x bundle, 50 Hz
%    R' = 25 m$\Omega$/km, X' = 0.30 $\Omega$/km, C' = 12.5 nF/km,
%    G' = 0, I_max = 2.0 kA, l = 400 km
% =============================================================
\providecommand{\acfig}[1]{design/ac-transmission/fig/#1}

\intermezzo{Natural Power}{Surge-impedance loading of an AC line}

% ---------------------------------------------------------------
\begin{frame}{Natural power and surge-impedance loading}
  \setlength{\abovedisplayskip}{5pt}\setlength{\belowdisplayskip}{5pt}
  \noindent{\small A line terminated in its own characteristic impedance
  $\underline{Z}_w$ carries its \emph{natural power}: the reactive power
  generated by the shunt capacitance exactly cancels the reactive power
  consumed by the series inductance, so the line exchanges no reactive power
  and its voltage magnitude stays constant. North American practice calls the
  same quantity the \emph{surge-impedance loading} (SIL).}

  \[ \underline{Z}_w
     = \sqrt{\frac{\underline{Z}'}{\underline{Y}'}}
     = \sqrt{\frac{R'+j\cdot\omega\cdot L'}{G'+j\cdot\omega\cdot C'}}
     \;\;\approx\;\;\sqrt{\frac{L'}{C'}} = Z_w
     \qquad
     P_{nat} = \mathrm{SIL} = \frac{U^2}{Z_w} \]

  {\small
  \begin{itemize}
    \setlength{\itemsep}{0.2em}
    \item On an overhead line $\omega\cdot L'\gg R'$ and
          $\omega\cdot C'\gg G'$, so the real $\sqrt{L'/C'}$ is within a
          percent of $\left|\underline{Z}_w\right|$ --- the value $P_{nat}$
          is quoted from.
    \item $Z_w$ is pure conductor geometry: independent of line length, and of
          frequency in the lossless form.
    \item $P_{nat}$ is a reference, not a rating: below it the line sources
          reactive power, above it absorbs it.
  \end{itemize}}
  \atnote{$U$ line-to-line voltage; $R',L',G',C'$ per-unit-length resistance,
  inductance, shunt conductance and capacitance; $\omega=2\cdot\pi\cdot f$.}
\end{frame}

% ---------------------------------------------------------------
\begin{frame}{Natural power --- definition}
  \begin{columns}[T,onlytextwidth]
    \begin{column}{0.53\textwidth}
      {\small Voltage and current at a distance $x$ from the receiving end
      follow from the two-port equations of the uniform line:}
      \[ \begin{aligned}
        \underline{U}(x) &= \underline{U}_R\cdot\cosh\!\left(\underline{\gamma}\cdot x\right)
                          + \underline{Z}_w\cdot\underline{I}_R\cdot\sinh\!\left(\underline{\gamma}\cdot x\right)\\[0.15em]
        \underline{I}(x) &= \frac{\underline{U}_R}{\underline{Z}_w}\cdot\sinh\!\left(\underline{\gamma}\cdot x\right)
                          + \underline{I}_R\cdot\cosh\!\left(\underline{\gamma}\cdot x\right)
      \end{aligned} \]
      {\small Closing the line with $\underline{Z}_w$ makes
      $\underline{I}_R=\underline{U}_R/\underline{Z}_w$, and both collapse to a
      single travelling wave --- no reflection:}
      \[ \underline{U}(x)=\underline{U}_R\cdot e^{\underline{\gamma}\cdot x},
         \quad
         \underline{I}(x)=\underline{I}_R\cdot e^{\underline{\gamma}\cdot x} \]
      {\small The magnitude then falls only by the attenuation
      $e^{-\alpha\cdot l}$, and $\underline{U}/\underline{I}=\underline{Z}_w$
      everywhere.}
    \end{column}
    \begin{column}{0.44\textwidth}\centering
      \topoimg[width=\linewidth,height=0.46\textheight,keepaspectratio]{\acfig{np_matched_profile.pdf}}\\[0.25em]
      {\color{atGray}\footnotesize Matched line, 400\,km --- only the
      complex $\underline{Z}_w$ matches perfectly.}
    \end{column}
  \end{columns}
\end{frame}

% ---------------------------------------------------------------
\begin{frame}{Natural power --- the reactive balance}
  \begin{columns}[T,onlytextwidth]
    \begin{column}{0.52\textwidth}
      {\small The same result follows from the reactive power balance of a line
      of length $l$ at line-to-line voltage $U$ and phase current $I$:}
      \[ \begin{aligned}
        Q_C &= \omega\cdot C'\cdot l\cdot U^2 && \text{(shunt, generated)}\\[0.15em]
        Q_L &= 3\cdot\omega\cdot L'\cdot l\cdot I^2 && \text{(series, consumed)}
      \end{aligned} \]
      {\small Natural loading is the point at which the two cancel. Setting
      $Q_C=Q_L$ eliminates $l$ and $\omega$ and gives
      $I=U/\!\left(\sqrt{3}\cdot Z_w\right)$, hence}
      \[ \boxed{\;P_{nat} = \sqrt{3}\cdot U\cdot I
                = \frac{U^2}{\sqrt{L'/C'}} = \frac{U^2}{Z_w}\;} \]
    \end{column}
    \begin{column}{0.45\textwidth}
      \renewcommand{\arraystretch}{1.15}
      {\footnotesize
      \begin{tabular}{@{}l l r@{}}
        \textbf{Quantity} & & \textbf{Value} \\
        \hline
        $R'$ & m$\Omega$/km & 25 \\
        $X'$ & $\Omega$/km  & 0.30 \\
        $L'=X'/\omega$ & mH/km & 0.955 \\
        $C'$ & nF/km & 12.5 \\
        \hline
        $Z_w=\sqrt{L'/C'}$ & $\Omega$ & 276.4 \\
        $\left|\underline{Z}_w\right|$ & $\Omega$ & 276.9 \\
        $\beta\cdot l$ & $^\circ$ & 24.88 \\
        \hline
        $P_{nat}$ & MW & 579 \\
        \hline
      \end{tabular}}\\[0.4em]
      {\color{atGray}\footnotesize 400\,kV, 2$\times$ bundle, 50\,Hz,
      $l=400$\,km, $G'=0$.}
    \end{column}
  \end{columns}
\end{frame}

% ---------------------------------------------------------------
\begin{frame}{Voltage profile along the line}
  \begin{columns}[c,onlytextwidth]
    \begin{column}{0.56\textwidth}\centering
      \topoimg[width=\linewidth,height=0.64\textheight,keepaspectratio]{\acfig{np_voltage_profile.pdf}}
    \end{column}
    \begin{column}{0.42\textwidth}
      {\small
      \begin{itemize}
        \setlength{\itemsep}{0.4em}
        \item Both terminals held at 1.0\,p.u.; the transmission angle
              $\vartheta$ is solved for each loading.
        \item Below $P_{nat}$ the shunt capacitance dominates and the middle
              of the line rises --- 1.024\,p.u.\ at no load.
        \item At $P_{nat}$ the profile is flat to four decimals.
        \item Above $P_{nat}$ the series reactance dominates and the middle
              sags --- 0.873\,p.u.\ at $2\cdot P_{nat}$, with
              $\vartheta=63^\circ$ across the line.
      \end{itemize}}
    \end{column}
  \end{columns}
\end{frame}

% ---------------------------------------------------------------
\begin{frame}{Reactive balance and the effect of the line losses}
  \begin{columns}[c,onlytextwidth]
    \begin{column}{0.57\textwidth}\centering
      \topoimg[width=\linewidth,height=0.56\textheight,keepaspectratio]{\acfig{np_reactive_balance.pdf}}
    \end{column}
    \begin{column}{0.41\textwidth}
      {\small The net reactive power absorbed by the line follows from the
      terminal powers,
      $Q_{line}=\Im\!\left\{\underline{S}_S\right\}-\Im\!\left\{\underline{S}_R\right\}$.
      \begin{itemize}
        \setlength{\itemsep}{0.35em}
        \item Lossless line: the zero crossing reproduces $U^2/Z_w$ to six
              decimals.
        \item Real line: 566\,MW, i.e.\ 2.2\,\% below the 579\,MW of
              $U^2/Z_w$.
        \item Cause --- the sending-end current carries the load \emph{plus}
              the losses, so $3\cdot I^2\cdot X'\cdot l$ is larger and the
              balance is reached at a lower transfer.
      \end{itemize}}
    \end{column}
  \end{columns}
  \atnote{Two-port solution; verified against a 20-segment pandapower model (deviation $<0.1$\,\%).}
\end{frame}

% ---------------------------------------------------------------
\begin{frame}{Natural power of typical line types}
  \centering
  \renewcommand{\arraystretch}{1.25}
  \begin{tabular}{@{}l r r r r r@{}}
    \textbf{Line} & \textbf{$U$} & \textbf{$C'$} & \textbf{$Z_w$}
                  & \textbf{$P_{nat}$} & \textbf{$S_{th}/P_{nat}$} \\
                  & kV & nF/km & $\Omega$ & MW & \\
    \hline
    230\,kV, 1$\times$ Drake     & 230 &  9.0 & 376.1 &  141 & 2.8 \\
    345\,kV, 2$\times$ bundle    & 345 & 11.8 & 292.3 &  407 & 2.6 \\
    \atrowhighlight
    400\,kV, 2$\times$ bundle    & 400 & 12.5 & 276.4 &  579 & 2.4 \\
    500\,kV, 3$\times$ bundle    & 500 & 13.5 & 256.6 &  974 & 2.3 \\
    765\,kV, 4$\times$ bundle    & 765 & 13.5 & 250.8 & 2334 & 2.0 \\
    \hline
  \end{tabular}

  \vspace{0.6em}
  {\small
  \begin{itemize}
    \setlength{\itemsep}{0.2em}
    \item $Z_w$ spans only 251--376\,$\Omega$ while $P_{nat}$ spans
          141--2334\,MW: $P_{nat}$ is a $U^2$ quantity.
    \item Bundling and tighter phase spacing are the only levers on $Z_w$ ---
          both lower $L'$ and raise $C'$.
    \item $S_{th}/P_{nat}$ decides where the thermal plateau of the
          loadability curve ends.
    \item AC cables sit at $Z_w\approx20\ldots40\,\Omega$: their $P_{nat}$ is
          far above the thermal rating, so charging current is the binding
          constraint instead.
  \end{itemize}}
\end{frame}

% ---------------------------------------------------------------
\begin{frame}{Loadability --- natural power is a reference, not a limit}
  \begin{columns}[c,onlytextwidth]
    \begin{column}{0.55\textwidth}\centering
      \topoimg[width=\linewidth,height=0.62\textheight,keepaspectratio]{\acfig{np_loadability.pdf}}
    \end{column}
    \begin{column}{0.43\textwidth}
      {\footnotesize Criteria: $I\le 2.0$\,kA, $U_R\ge 0.95$\,p.u.,
      $\vartheta\le 30^\circ$, no compensation.}\\[0.45em]
      \renewcommand{\arraystretch}{1.15}
      {\footnotesize
      \begin{tabular}{@{}l l@{}}
        \multicolumn{2}{@{}l}{\textbf{(a) load at receiving end}}\\
        \hline
        $\le 80$\,km   & thermal, $2.3\ldots2.4\cdot P_{nat}$ \\
        120--400\,km   & voltage drop \\
        $\ge 500$\,km  & angular stability \\
        382\,km        & falls through $1\cdot P_{nat}$ \\
        \hline
      \end{tabular}}\\[0.45em]
      {\footnotesize
      \begin{tabular}{@{}l l@{}}
        \multicolumn{2}{@{}l}{\textbf{(b) two stiff systems}}\\
        \hline
        $\le 160$\,km  & thermal \\
        $\ge 200$\,km  & angular stability \\
        472\,km        & falls through $1\cdot P_{nat}$ \\
        \hline
      \end{tabular}}
    \end{column}
  \end{columns}
  \atnote{The boundary condition matters: holding $U_R$ at 1.0\,p.u.\ (case b)
  removes the voltage-drop criterion entirely and shifts the $1\cdot P_{nat}$
  crossing by 90\,km. Case (a) is the classical St.\ Clair configuration.}
\end{frame}

% ---------------------------------------------------------------
\begin{frame}{Light load --- the Ferranti effect}
  \begin{columns}[c,onlytextwidth]
    \begin{column}{0.55\textwidth}\centering
      \topoimg[width=\linewidth,height=0.44\textheight,keepaspectratio]{\acfig{np_ferranti.pdf}}
    \end{column}
    \begin{column}{0.43\textwidth}
      {\small Open-ended line, $\underline{I}_R=0$:
      \[ \frac{\underline{U}_R}{\underline{U}_S}
         = \frac{1}{\cosh\!\left(\underline{\gamma}\cdot l\right)}
         \;\;\longrightarrow\;\;
         \frac{1}{\cos\!\left(\beta\cdot l\right)} \]
      \begin{itemize}
        \setlength{\itemsep}{0.3em}
        \item $\beta\cdot l=24.88^\circ$ gives 1.102\,p.u.\ at the open end,
              from 251\,MVAr of charging power.
        \item A 60\,\% shunt reactor brings it back to 0.984\,p.u.; full
              compensation gives 0.918\,p.u.
        \item The degree follows from the light-load voltage, but it moves
              $P_{nat}$ too.
      \end{itemize}}
    \end{column}
  \end{columns}
\end{frame}

% ---------------------------------------------------------------
\begin{frame}{Compensation moves the natural power}
  \begin{columns}[c,onlytextwidth]
    \begin{column}{0.52\textwidth}\centering
      \topoimg[width=\linewidth,height=0.56\textheight,keepaspectratio]{\acfig{np_compensation.pdf}}
    \end{column}
    \begin{column}{0.46\textwidth}
      {\small For \emph{uniformly distributed} compensation the scaling is
      exact:
      \[ \begin{aligned}
        X'&\to X'\cdot\left(1-k\right)
          &&\Rightarrow\;\; \frac{P_{nat}}{\sqrt{1-k}}\\[0.3em]
        C'&\to C'\cdot\left(1-m\right)
          &&\Rightarrow\;\; P_{nat}\cdot\sqrt{1-m}
      \end{aligned} \]}
      {\small
      \begin{itemize}
        \setlength{\itemsep}{0.3em}
        \item $k=60$\,\% raises $P_{nat}$ by $1.581\times$; $m=60$\,\% lowers
              it by $0.632\times$.
        \item Real banks are lumped at substations, so the laws are a
              first-order guide only: a lumped 30\,\% series bank at
              $1.4\cdot P_{nat}$ reduces $\vartheta$ from $37.6^\circ$ to
              $25.6^\circ$ and turns the sending-end reactive flow from
              $+89$ to $-30$\,MVAr.
      \end{itemize}}
    \end{column}
  \end{columns}
\end{frame}
